\section{Introduction}\label{p:introduction}

A linear projection can concentrate a measure even when the original
measure has no atoms. The amount of concentration depends on how much mass
can lie near the affine subspaces that a projection collapses. This paper
studies that dependence for random exact projections and uses it to analyze
what a finite memory state can retain from noiseless Gaussian equations.

We consider three questions. The first asks for an integrated moment of a
projected density. Let \(\nu\) be a finite positive measure in \(\R^d\)
whose mass on every ball of radius \(r\) is at most \(Kr^\beta\), and let
\(P\) be a random matrix with \(k\) orthonormal rows. When \(P_\#\nu\)
has density \(g_P\), an integrated estimate controls the average over
\(P\) of \(\int g_P^q\). Such an estimate measures concentration throughout
the label space. Theorem~\ref{p:frame-moment} proves a bound with explicit
dependence on the dimension, support radius, and local-mass exponent.

The second question asks for an actual positive dominating measure after a
finite routing decision. Let \(\sigma\) be uniform probability on
\(S^{d-1}\), and let \(\mu_B\) be \(\sigma\) conditioned on a spherical
cap \(B\). A routing rule observes a Gaussian matrix \(X\) and the exact
label \(Xs\), and selects an output \(w\) with probability
\(K_w(X,Xs)\). Its average likelihood \(f_w(s)=\E_XK_w(X,Xs)\)
weights the parent measure to \(f_w\mu_B\). We seek a positive sum of
normalized cap measures that dominates this measure, with a controlled
cost for small child radii. The positive-domination theorem in
Section~\ref{p:cap-section} constructs a countable such sum and bounds
its cost by a power of a reference routing probability.

The third question concerns evaluation at a specified label. For a
measurable function \(h:S^{d-1}\to[0,1]\), the ambient Riesz functional
\[
 V(h)=\sup_{z\in\R^d}\int h(s)\norm{s-z}^{-d/2}\,d\sigma(s)
\]
is the supremum of the raw Riesz potentials of \(h\sigma\) over ambient
centers, and it controls mass near every center. If \(h\sigma\) is restricted
to a ball about \(z\), Section~\ref{p:riesz-section} estimates a Gaussian projected
density at \(Xz+b\) for every fixed offset \(b\). It uses one explicitly
chosen density version for all offsets. This qualification matters:
an integrated density norm is unchanged by modifications on a
matrix-dependent graph, so it cannot by itself assign or bound the values
on that graph.

All three estimates give a precision bound for finite-state regression.
The signal is uniform on \(S^{d-1}\), and each fresh observation is a
standard Gaussian feature together with its exact inner product with the
signal. Write \(T\) for the learner's deterministic sample horizon.
In the bit-memory model of Section~\ref{p:geometry-model}, a learner
with \(M(d)=o(d^2)\) retained bits that achieves angular error at most
\(\epsilon(d)\), where \(0<\epsilon(d)\le1/10\), with average probability
at least \(2/3\)
requires
\[
 T(d)\ge c\,d\log(1/\epsilon(d))
\]
for an absolute \(c>0\) and all sufficiently large \(d\), as stated in
Theorem~\ref{p:regression}. The dimension threshold may depend on the
memory sequence. Under the same accuracy and success requirements, the
integrated projection argument also gives the explicit tradeoff
\[
 T\ge \frac{c\,d\log(1/\epsilon)}{1+M/d^2}
\]
for every nonnegative integer \(M\) and all \(d\) above an absolute
threshold (Corollary~\ref{p:finite-memory-tradeoff}). In particular, the
same precision bound holds with absolute constants throughout \(M\le d^2\).
The underlying success-probability bounds
\eqref{p:haar-success}, \eqref{p:cap-fixed-stop-explicit}, and
Proposition~\ref{p:riesz-success} retain explicit dependence on arbitrary
\(M\).

\subsection{Projection geometry and earlier methods}

The local growth condition belongs to classical potential theory.
Mattila's work on Hausdorff dimension and orthogonal projections uses
inverse-distance energy averaging and compares positive-dimensional
Hausdorff measure with the existence of a measure satisfying a ball-growth
bound~\cite[Lemma~5.1, p.~236, and Theorem~5.4, p.~237]{Mattila1975}.
Mattila describes the potential-theoretic projection method as following
Kaufman's work~\cite[Introduction, p.~227]{Mattila1975}. Here the growth
measure is supplied by a hypothesis or constructed from a cap test. The quantitative estimates require successive
distances to affine spans and moments whose order grows with \(d\).

Affine simplex volumes also have an established role in integral geometry.
Drury's affine-plane change of variables contains a simplex-volume
Jacobian~\cite[Lemma~1, pp.~497--498]{Drury1984}. In our frame calculation,
the product of successive affine heights is the volume of the associated
parallelepiped, and the \(k\) projected rows raise its inverse to the
\(k\)th power. We derive that row-Jacobian and its probability normalization
directly. The exact density of a rectangular block of a Haar orthogonal
matrix is classical; see Jiang~\cite[Lemma~2.5]{Jiang2006}.
Our derivation keeps its normalization explicit before the affine-height
change of variables. Classical beta descriptions of a uniform direction projected
onto a subspace are discussed by Frankl and
Maehara~\cite{FranklMaehara1990}; Lemma~\ref{p:coordinate-density}
gives the vector density and conditional direction law used here.

The regression application sits within a broader study of learning with
limited retained information. Shamir's finite-message protocols include a
bounded-memory online subclass~\cite[Section~2]{Shamir2014}.
Steinhardt, Valiant, and Wager related memory, communication, and statistical
queries~\cite[Section~1.1]{SVW2016}; Raz subsequently proved a qualitative
quadratic-memory versus exponential-sample separation for parity
learning~\cite[Theorem~1]{Raz2016}. For regression, Steinhardt and Duchi
established memory-dependent minimax risk bounds in a sparse noisy
model~\cite[Section~1.1]{SteinhardtDuchi2015}.

Sharan, Sidford, and Valiant developed a nonuniform finite-width
branching-program analysis for continuous regression. Their principal
theorem concerns independent isotropic Gaussian features with independent
uniform additive noise of half-width \(2^{-d/5}\). With at most \(d^2/4\)
bits of memory, success probability at least \(2/3\) for a uniformly
distributed unit signal, and Euclidean accuracy
\(d^{-r}\) in its stated range \(r\le O(d/\log d)\), it gives an
\(\Omega(d\log r)\) sample bound~\cite[Theorem~1]{SSV2019}.
Their projection argument expands high moments into independent copies
and uses successive orthogonalization~\cite[Section~7, Lemma~12]{SSV2019}.
That is a close methodological predecessor of our tuple calculations.
Their lemma treats one noisy interval under a global \(L^2\) density
condition. The estimates here treat exact multirow labels under local
ball or Riesz-potential conditions. Dagan, Kur, and Shamir obtained related
quadratic-space obstructions for approximate null vectors of Gaussian
rows and for empirical residual error in a different equation
model~\cite[Theorems~3 and~10]{DKS2019}. These comparisons concern
distinct statistical targets and observation laws.

\subsection{How the three arguments work}

For the integrated estimate, expand the \(q\)th density moment using
\(q\) points of the measure. The collision of their labels is controlled
by the projection of the \(q-1\) differences from the first point.
Successive coordinates of an orthonormal frame have ellipsoidal densities;
their determinant powers telescope, leaving one nonnegative determinant
power and an explicit \(\sqrt d\) scale. Gram--Schmidt converts the
remaining change of variables into inverse affine heights. A tube about
each preceding affine span is covered by balls, so the local mass
hypothesis integrates these heights in reverse order. This proves
Theorem~\ref{p:frame-moment}.

The first application separates a Gaussian matrix into its positive radial
factor and an independent orthonormal frame. The routing rule may use the
whole radial factor; averaging that dependence produces weights on frame
and projected label. The projected support occupies a \(k\)-ball.
Its volume cancels the dimensional scale in the density norm and returns
the same radius exponent as the incoming local-mass bound. The backward
argument therefore preserves the ambient dimension through every block.
Section~\ref{p:projection} completes this route with local stopping-index
counting and the common full-data bound of Section~\ref{p:geometry-model}.

For positive cap domination, the order is different. Bounded nonnegative
tests on a countable cap dictionary imply a local ball bound for a rescaled
test measure. The normalized instance of the integrated theorem bounds the
test of \(f_w\mu_B\). To turn all these inequalities into a positive
dominating measure, the proof uses a weighted sum of cap densities.
Finitely many caps occur at each scale and the weighted cost makes the
fine-scale \(L^1\) tail vanish. This gives norm compactness. Strict
separation of the closed downward convex set then produces exactly the
bounded tests already controlled. The resulting positive sums propagate
the unnormalized joint laws of signal and state, with a separate reference
probability for each parent cap.

The specified-label argument uses \(q\) vectors anchored at \(z\),
rather than \(q-1\) differences anchored at a sampled point. Their images
under the Gaussian rows have covariance equal to the anchored Gram matrix.
The resulting ball-average estimate is uniform in every fixed offset.
Taking a single lower-limit version and applying Fatou's lemma on each
offset graph preserves that estimate for exact densities. Operator-norm
bins and radial shells then control the Riesz potential of earlier
continuations. This route retains a factor of order
\(1+\log(1/\epsilon)\) from its shell count and a global terminal clock;
its own short-run argument controls both.

The three analytic conclusions and the costs of their regression
applications are compared below. Here \(T\) is the deterministic sample
horizon, and a stopping index is charged only where its value enters the
corresponding argument.

\begin{center}
\small
\begin{tabular}{@{}>{\raggedright\arraybackslash}p{.18\textwidth}
                  >{\raggedright\arraybackslash}p{.37\textwidth}
                  >{\raggedright\arraybackslash}p{.39\textwidth}@{}}
\toprule
Estimate & Analytic guarantee & Regression argument \\
\midrule
Integrated projection & A density moment integrated over the frame and
label space & Backward local-mass propagation; stopping indices counted
within each block \\[3pt]
Positive cap domination & An actual dominating sum of cap measures, with
weighted radius cost & Forward propagation of joint signal--state laws;
one final sum over stopping indices \\[3pt]
Fixed-offset density & One density version with a moment bound at
every fixed affine offset & Backward Riesz-potential propagation; global
terminal tags and a logarithmic shell count \\
\bottomrule
\end{tabular}
\end{center}

Section~\ref{p:geometry-model} supplies the probability normalizations,
finite-state model, Borel-kernel reduction, and common full-data endpoint.
Section~\ref{p:projection} proves the integrated estimates and their
backward application; these two sections suffice for the regression
consequence and the explicit memory tradeoff. Section~\ref{p:cap-section}
uses the normalized projection moment to construct positive dominating
measures and propagate them forward. Section~\ref{p:riesz-section} first
proves its specified-offset estimate, then develops the Riesz recurrence
and its parameter argument. The cap dual constraints test averages on
spherical caps; they are distinct from integration on affine subspheres.
